% ============================================================================
% Protocolo de Pruebas - Prototipo 3
% Dashboard de rendimiento, gráfica de tendencias, reporte de desviaciones
% y pruebas de integridad y usuario.
%
% Documento autónomo: puede compilarse directamente con pdfLaTeX.
% Los resultados de las columnas "Aprobado" y "Observaciones" provienen de la
% ejecución real de la suite tests/e2e (ver sección "Ejecución de las pruebas").
% ============================================================================
\documentclass[12pt,letterpaper]{article}

\usepackage[utf8]{inputenc}
\usepackage[T1]{fontenc}
\usepackage[spanish,es-noindentfirst]{babel}
\usepackage[margin=2.5cm]{geometry}
\usepackage{array}
\usepackage{tabularx}
\usepackage{longtable}
\usepackage{float}
\usepackage{xcolor}
\usepackage{hyperref}

% Tipo de columna usado en las tablas del protocolo.
\newcolumntype{L}[1]{>{\raggedright\arraybackslash}p{#1}}

% Compatibilidad con el comando del documento principal.
\newcommand{\ttiiTOC}[1]{\texttt{#1}}

% Nombre de método de prueba (permite guiones bajos sin escaparlos).
\newcommand{\mth}[1]{\texttt{\detokenize{#1}}}

\title{\textbf{Protocolo de Pruebas\\[0.3em]Prototipo 3}}
\author{AI Presentation Analyzer}
\date{\today}

\begin{document}
\maketitle

\section*{Protocolo de Pruebas del Prototipo 3}
\addcontentsline{toc}{section}{Protocolo de Pruebas del Prototipo 3}

Para facilitar la trazabilidad de las pruebas del prototipo 3, se utiliza un prefijo de tres letras y un número consecutivo. La Tabla \ref{tab:nomenclatura-pruebas-prototipo-3} presenta las categorías empleadas.

\begin{table}[H]
\centering
\renewcommand{\arraystretch}{1.2}
\begin{tabularx}{\textwidth}{|>{\centering\arraybackslash}p{1.5cm}|p{4.6cm}|X|}
\hline
\textbf{Sigla} & \textbf{Módulo o categoría} & \textbf{Significado dentro del protocolo} \\ \hline
PRF & Dashboard de rendimiento & Entrada a la vista, recuperación de los JSON y construcción de las tarjetas de resumen. \\ \hline
GRF & Gráfica de tendencias & Series de tiempo, \textit{tooltip}, desplazamiento horizontal y pestañas de grupo. \\ \hline
DVS & Reporte de desviaciones & Apertura del modal, cálculo de porcentajes y clasificación de diapositivas. \\ \hline
INT & Integridad & Consistencia entre el análisis, el reporte de tiempos y la base de datos. \\ \hline
USR & Pruebas de usuario & Escenarios derivados de la interacción directa con el dashboard. \\ \hline
\end{tabularx}
\caption{Nomenclatura de los identificadores de las pruebas del prototipo 3.}
\label{tab:nomenclatura-pruebas-prototipo-3}
\end{table}

% ===========================================================================
\section*{Ejecución de las pruebas}
\addcontentsline{toc}{section}{Ejecución de las pruebas}

Todas las pruebas de este protocolo están automatizadas como pruebas de extremo a extremo (E2E) que arrancan la interfaz real (CustomTkinter) y disparan los callbacks productivos. Se ejecutan desde la raíz del repositorio con:

\begin{verbatim}
python -m tests.e2e.run_e2e
\end{verbatim}

El ejecutor guarda la salida cruda y un resumen en \texttt{docs/prototipo\_final/evidencias/<fecha\_hora>/}. La corrida documentada en este protocolo corresponde a la carpeta:

\begin{center}
\texttt{docs/prototipo\_final/evidencias/20261018\_233448/}
\end{center}

\begin{table}[H]
\centering
\renewcommand{\arraystretch}{1.2}
\begin{tabularx}{\textwidth}{|p{4.5cm}|X|}
\hline
\textbf{Dato} & \textbf{Valor} \\ \hline
Fecha de ejecución & 2026-10-18 23:38:33 \\ \hline
Intérprete & Python 3.14.6 (64 bits) \\ \hline
Plataforma & Windows 11 (10.0.26300) \\ \hline
Interfaz & CustomTkinter / Tk 3.14.6 (con display) \\ \hline
SQLite & 3.50.4 (stdlib) \\ \hline
\end{tabularx}
\caption{Entorno de ejecución de las pruebas del prototipo 3.}
\label{tab:entorno-ejecucion-tt3}
\end{table}

\begin{table}[H]
\centering
\renewcommand{\arraystretch}{1.2}
\begin{tabularx}{\textwidth}{|X|>{\centering\arraybackslash}p{2.6cm}|}
\hline
\textbf{Módulo de la suite} & \textbf{Resultado} \\ \hline
\texttt{tests.e2e.test\_e2e\_negocio} & 15 pruebas, 0 fallos \\ \hline
\texttt{tests.e2e.test\_e2e\_gui} & 15 pruebas, 0 fallos \\ \hline
\texttt{tests.e2e.test\_e2e\_rendimiento\_gui} & 67 pruebas, 0 fallos \\ \hline
\textbf{Total} & \textbf{97 pruebas, 0 fallos, 0 errores, resultado OK} \\ \hline
\end{tabularx}
\caption{Resultado global de la corrida de evidencias.}
\label{tab:resultado-global-tt3}
\end{table}

Las 67 pruebas del módulo \texttt{test\_e2e\_rendimiento\_gui} son la evidencia directa de las 67 filas de este protocolo (la Tabla \ref{tab:trazabilidad-tt3} mapea cada identificador con su método de prueba).

% ===========================================================================
\clearpage
\textbf{Módulo: Dashboard de Rendimiento}

\textbf{Objetivo:} Validar la recuperación del análisis de la IA y del reporte de tiempos, así como la construcción de las tarjetas de resumen y la gráfica de tendencias.

\textbf{Descripción del flujo:}
\begin{enumerate}
\item Acceder mediante «Ver rendimiento», disponible solo si existe al menos una sesión registrada.
\item Configurar la interfaz en modo secundario.
\item Resolver el identificador de la materia.
\item Recuperar los JSON de análisis y tiempos.
\item Construir las cinco tarjetas de resumen y la gráfica.
\item Restaurar la vista mediante «Volver a inicio».
\end{enumerate}

\textbf{Arquitectura del módulo:}
\begin{verbatim}
Dashboard de Rendimiento
|-- views/performance_gui.py
|-- views/navigator.py
|-- widgets/performance_mode/
|   |-- performance_summary_panel.py
|   |-- performance_chart_panel.py
|   +-- performance_deviations_panel.py
|-- controller/
|   |-- presentation_controller.py
|   +-- subject_controller.py
+-- data/queries.py, database_manager.py
\end{verbatim}

\textbf{Componentes que probar:}
\begin{table}[H]
\centering
\begin{tabular}{|L{0.34\textwidth}|L{0.18\textwidth}|L{0.38\textwidth}|}
\hline
\textbf{Componente} & \textbf{Tipo} & \textbf{Descripción} \\ \hline
Vista de rendimiento & UI/Frontend & Orquesta la carga de los JSON y monta los paneles. \\ \hline
Tarjetas de resumen & UI/Frontend & Muestran calificación, tiempos y desviación. \\ \hline
Controladores & Backend & Obtienen el análisis, el reporte y el identificador de la materia. \\ \hline
Consultas & Base de datos & Extraen los datos desde \texttt{Historial\_de\_Versiones}. \\ \hline
\end{tabular}
\caption{Componentes del módulo de dashboard de rendimiento.}
\label{tab:componentes-modulo-dashboard}
\end{table}

\textbf{Cobertura de pruebas --- Dashboard de Rendimiento:}
{\tiny
\renewcommand{\arraystretch}{1.15}
\begin{longtable}{|L{1.1cm}|L{2.3cm}|L{3.2cm}|L{3.6cm}|L{1.2cm}|L{3.6cm}|}
\hline
\textbf{ID} & \textbf{Caso} & \textbf{Descripción} & \textbf{Resultado esperado} & \textbf{Aprobado} & \textbf{Observaciones} \\ \hline
\endhead
PRF-01 & Entrada a la vista & Seleccionar «Ver rendimiento». & Se despliega el dashboard. & Sí & Ejecutado en \mth{test_prf01_entrada_a_la_vista}: la vista queda activa y se montan el contenedor, las tarjetas de resumen y la gráfica. \\ \hline
PRF-02 & Carga de análisis & Recuperar la evaluación generada durante el análisis. & Se recupera el JSON de análisis. & Sí & \mth{test_prf02_carga_de_analisis} comprobó que el índice de diapositivas coincide con el JSON y que \texttt{total\_diapositivas} se conserva. \\ \hline
PRF-03 & Carga de tiempos & Abrir una presentación con sesiones. & Se recupera el reporte de tiempos. & Sí & \mth{test_prf03_carga_de_tiempos} verificó que el grupo \texttt{3CM2} aparece como pestaña con su duración real acumulada. \\ \hline
PRF-04 & Resumen con datos & Revisar las cinco tarjetas. & Se muestran todos los indicadores. & Sí & \mth{test_prf04_resumen_con_datos} leyó los valores reales: «Planeó docente» = 60m 00s, «Duración real» = 7m 00s, subtítulo «4/4 diapositivas» y «Desviación» = 0\,\%. \\ \hline
PRF-05 & Resumen sin grupos & Abrir una presentación sin sesiones. & Se abre la interfaz de rendimiento. & No & Por diseño: \texttt{\_ver\_rendimiento\_ui} exige \texttt{historial\_sesiones} y muestra el modal «Sin datos de rendimiento» (\mth{test_uhp02_ver_rendimiento_sin_sesiones_muestra_modal}). A nivel de panel, \mth{test_prf05_resumen_sin_grupos} confirmó que con grupos vacío el resumen queda en «—» y «sin grupo». \\ \hline
PRF-06 & Sin análisis & Abrir una presentación no analizada. & Se informa la ausencia de datos. & No & El control de acceso también exige el reporte de tiempos, por lo que no se navega. Si faltara el análisis, \mth{test_prf06_sin_analisis} comprobó que la gráfica dibuja «Sin datos de análisis» y el contador queda en «0/0 diapositivas». \\ \hline
PRF-07 & Tiempo IA & Mostrar el tiempo de exposición recomendado. & El total coincide con la suma. & Sí & \mth{test_prf07_tiempo_ia_es_la_suma}: con el campo total ausente, el dashboard sumó las diapositivas (60+120+90 = 270 s) y mostró «4m 30s». \\ \hline
PRF-08 & Grupo inicial & Abrir una presentación con dos grupos. & El dashboard selecciona el primer grupo. & Sí & \mth{test_prf08_grupo_inicial}: con \texttt{\{3CM2, 3CM4\}} quedó activo \texttt{3CM2}, el primero del reporte. \\ \hline
PRF-09 & Análisis corrupto & Introducir un JSON inválido. & Se controla la excepción. & Sí & \mth{test_prf09_analisis_corrupto}: \texttt{json.loads} falló dentro del \texttt{try/except}, el análisis quedó vacío y la vista continuó funcionando sin excepción. \\ \hline
PRF-10 & Reporte corrupto & Introducir un reporte inválido. & Los grupos quedan vacíos. & Sí & \mth{test_prf10_reporte_corrupto}: el reporte inválido dejó \texttt{grupos = \{\}} sin lanzar error. \\ \hline
PRF-11 & Regreso al inicio & Seleccionar «Volver a inicio». & Se restaura la vista principal. & Sí & \mth{test_prf11_regreso_al_inicio}: la vista volvió a «main» y las referencias del contenedor y de la gráfica quedaron en \texttt{None}. \\ \hline
PRF-12 & Eliminación durante la vista & Intentar eliminar la presentación. & No aplica por diseño. & N/A & \mth{test_prf12_eliminacion_no_aplica} confirmó que la vista no expone acciones de eliminación. \\ \hline
PRF-13 & Calificación por etiqueta & Revisar la tarjeta de calificación. & Se muestra una etiqueta cualitativa. & Sí & \mth{test_prf13_calificacion_por_etiqueta}: con \texttt{zona\_global = "bueno"} la tarjeta mostró «Bueno» (etiqueta capitalizada, no un número). \\ \hline
PRF-14 & Desviación global & Comparar la duración real con la IA. & La tarjeta usa el color según el umbral. & Sí & \mth{test_prf14_desviacion_global_por_umbral} verificó «+200\,\%» (naranja), «-60\,\%» (ámbar) y el color verde dentro de $\pm 20$\,\%. \\ \hline
PRF-15 & Diapositivas presentadas & Observar el subtítulo de «Duración real». & Indica cuántas diapositivas tienen tiempo real. & Sí & \mth{test_prf15_diapositivas_presentadas}: con 2 de 4 diapositivas presentadas el subtítulo mostró «2/4 diapositivas». \\ \hline
PRF-16 & Tiempo planeado del grupo & Guardar dos sesiones y abrir el dashboard. & Usa el tiempo de la última sesión. & Sí & \mth{test_prf16_planeado_usa_la_ultima_sesion}: con sesiones de 1800 s y 2700 s mostró «45m 00s» (la última). \\ \hline
PRF-17 & Reapertura del dashboard & Salir del dashboard y volver a entrar. & Se reconstruye sin residuos. & Sí & \mth{test_prf17_reapertura_sin_residuos}: tras salir y volver a entrar, los grupos se recargaron y la vista quedó operativa. \\ \hline
\caption{Casos de prueba del módulo de dashboard de rendimiento.}
\label{tab:casos-prueba-dashboard}
\end{longtable}
}

\clearpage
\textbf{Módulo: Gráfica de Tendencias}

\textbf{Objetivo:} Validar el dibujado de las series de tiempo real y recomendado por la IA, el \textit{tooltip}, el desplazamiento horizontal y la actualización al cambiar de grupo.

\textbf{Descripción del flujo:}
\begin{enumerate}
\item Construir el lienzo, los ejes y la barra de desplazamiento.
\item Dibujar el área real y ambas series.
\item Mostrar el detalle de cada punto al pasar el cursor.
\item Cambiar de grupo y redibujar la gráfica y el resumen.
\item Ajustar el contenido al redimensionar la ventana.
\end{enumerate}

\textbf{Arquitectura del módulo:}
\begin{verbatim}
Gráfica de Tendencias
+-- widgets/performance_mode/performance_chart_panel.py
\end{verbatim}

\textbf{Componentes que probar:}
\begin{table}[H]
\centering
\begin{tabular}{|L{0.34\textwidth}|L{0.18\textwidth}|L{0.38\textwidth}|}
\hline
\textbf{Componente} & \textbf{Tipo} & \textbf{Descripción} \\ \hline
Lienzo de la gráfica & UI/Frontend & Dibuja ejes, series, área, puntos y detalles interactivos. \\ \hline
Pestañas de grupo & UI/Frontend & Cambian el grupo activo y notifican al resumen. \\ \hline
Barra de desplazamiento & UI/Frontend & Permite consultar presentaciones extensas. \\ \hline
\textit{Tooltip} & UI/Frontend & Muestra diapositiva, tiempo real y tiempo recomendado. \\ \hline
\end{tabular}
\caption{Componentes del módulo de gráfica de tendencias.}
\label{tab:componentes-modulo-grafica}
\end{table}

\textbf{Cobertura de pruebas --- Gráfica de Tendencias:}
{\tiny
\renewcommand{\arraystretch}{1.15}
\begin{longtable}{|L{1.1cm}|L{2.3cm}|L{3.2cm}|L{3.6cm}|L{1.2cm}|L{3.6cm}|}
\hline
\textbf{ID} & \textbf{Caso} & \textbf{Descripción} & \textbf{Resultado esperado} & \textbf{Aprobado} & \textbf{Observaciones} \\ \hline
\endhead
GRF-01 & Dibujado inicial & Abrir un grupo válido. & Se dibujan ejes, guías y series. & Sí & \mth{test_grf01_dibujado_inicial}: el lienzo contiene líneas (guías, ejes y series), el polígono del área, óvalos y las etiquetas 1 a 5 del eje X. \\ \hline
GRF-02 & Área real & Revisar el polígono bajo la curva real. & El área utiliza el color definido. & Sí & \mth{test_grf02_area_real}: se dibuja exactamente un polígono con relleno \#F4E2E7. \\ \hline
GRF-03 & Serie IA & Revisar la línea recomendada. & Se muestra una línea dorada discontinua. & Sí & \mth{test_grf03_serie_ia_discontinua}: se encontró la serie con color \#BC955C y trazo discontinuo. \\ \hline
GRF-04 & Tooltip & Pasar el cursor sobre un punto. & Se muestra la guía y el detalle. & Sí & \mth{test_grf04_tooltip_muestra_detalle}: al posar el cursor se crearon los elementos «hover» y el texto «Diapositiva N». \\ \hline
GRF-05 & Cursor fuera de rango & Alejar el cursor del punto. & El detalle desaparece. & Sí & \mth{test_grf05_cursor_lejos_del_punto}: con puntos separados más de 40 px, al alejar el cursor se eliminó todo el grupo «hover» (sin elementos). \\ \hline
GRF-06 & Desplazamiento & Usar más de doce diapositivas. & La barra permite consultar todas. & Sí & \mth{test_grf06_desplazamiento_entre_ventanas}: con 30 diapositivas \texttt{por\_vista < total} y al mover el inicio cambió la ventana visible. \\ \hline
GRF-07 & Cambio de grupo & Seleccionar otra pestaña. & La gráfica usa los nuevos datos. & Sí & \mth{test_grf07_cambio_de_grupo_redibuja}: al elegir \texttt{3CM4} cambió el grupo activo y los valores reales de los puntos. \\ \hline
GRF-08 & Sincronización & Cambiar de grupo. & El resumen también se actualiza. & Sí & \mth{test_grf08_sincronizacion_con_el_resumen}: el callback \texttt{on\_grupo\_cambiado} recibió el grupo seleccionado. \\ \hline
GRF-09 & Escala vertical & Usar un grupo con valores mayores. & La escala se ajusta al máximo. & Sí & \mth{test_grf09_escala_vertical}: el máximo vertical fue 200 (grupo) frente a 150 (IA), y 500 cuando la IA era mayor. \\ \hline
GRF-10 & Sin grupos & Abrir una presentación sin sesiones. & Se muestra el estado vacío. & N/A & El flujo no permite abrir el dashboard sin sesiones. A nivel de panel, \mth{test_grf10_sin_grupos} validó el estado vacío: pestañas «Sin grupos» y lienzo «Sin datos de análisis». \\ \hline
GRF-11 & Redimensionado & Cambiar el tamaño de la ventana. & La gráfica se adapta al nuevo ancho. & Sí & \mth{test_grf11_redimensionado}: el evento \texttt{<Configure>} reprogramó el dibujado (\texttt{\_resize\_id} no nulo). \\ \hline
GRF-12 & Muchas diapositivas & Usar una presentación extensa. & La interacción permanece fluida. & Sí & \mth{test_grf12_muchas_diapositivas}: con 120 diapositivas solo se pintan \texttt{por\_vista} puntos y no hubo errores. \\ \hline
GRF-13 & No presentada & Consultar una diapositiva sin tiempo real. & Se marca como no presentada. & Sí & \mth{test_grf13_diapositiva_no_presentada}: el detalle mostró «(no presentada)» para la diapositiva con tiempo real 0. \\ \hline
GRF-14 & Tiempo cero & Usar tiempo recomendado igual a cero. & El punto se dibuja sin errores. & Sí & \mth{test_grf14_tiempo_ia_cero_no_rompe}: con IA = 0 la escala se mantuvo positiva y los puntos se dibujaron sin división entre cero. \\ \hline
GRF-15 & Etiquetas de valor & Revisar las etiquetas sobre los puntos. & Se muestran los segundos reales. & Sí & \mth{test_grf15_etiquetas_de_valor}: con \texttt{por\_vista <= 12} se rotularon los valores reales (por ejemplo «20s»). \\ \hline
GRF-16 & Leyenda & Revisar la parte superior de la gráfica. & Se distingue real de recomendado. & Sí & \mth{test_grf16_leyenda}: se encontraron los textos «Tiempo real», «Recomendado IA» y el botón «Desviaciones». \\ \hline
GRF-17 & Límites del desplazamiento & Arrastrar la barra a los extremos. & No se puede salir del rango. & Sí & \mth{test_grf17_limites_del_desplazamiento}: con inicio 9999 se recortó a \texttt{total - por\_vista}, y con -50 a 0. \\ \hline
GRF-18 & Presentación corta & Abrir una presentación con menos de cinco diapositivas. & La gráfica se dibuja sin errores. & Sí & \mth{test_grf18_presentacion_corta}: con 2 diapositivas se trazaron 2 puntos y se rotuló el eje X sin error. \\ \hline
\caption{Cobertura de pruebas del módulo de gráfica de tendencias.}
\label{tab:cobertura-grafica}
\end{longtable}
}

\clearpage
\textbf{Módulo: Reporte de Desviaciones}

\textbf{Objetivo:} Validar la clasificación de diapositivas con duración insuficiente o excesiva y la presentación de hasta tres casos por categoría.

\textbf{Descripción del flujo:}
\begin{enumerate}
\item Abrir el modal mediante «Desviaciones».
\item Filtrar las diapositivas presentadas y calcular su porcentaje.
\item Clasificar, ordenar y seleccionar los tres casos principales.
\item Mostrar las dos columnas y cerrar mediante «Entendido» o la X.
\end{enumerate}

\textbf{Arquitectura del módulo:}
\begin{verbatim}
Reporte de Desviaciones
+-- widgets/performance_mode/performance_deviations_panel.py
\end{verbatim}

\textbf{Componentes que probar:}
\begin{table}[H]
\centering
\begin{tabular}{|L{0.34\textwidth}|L{0.18\textwidth}|L{0.38\textwidth}|}
\hline
\textbf{Componente} & \textbf{Tipo} & \textbf{Descripción} \\ \hline
Modal de desviaciones & UI/Frontend & Muestra dos columnas y hasta tres tarjetas por categoría. \\ \hline
Cálculo porcentual & Lógica & Compara el tiempo real con el recomendado. \\ \hline
Clasificación & Lógica & Ordena y limita los resultados por categoría. \\ \hline
\end{tabular}
\caption{Componentes del módulo de reporte de desviaciones.}
\label{tab:componentes-modulo-desviaciones}
\end{table}

\textbf{Cobertura de pruebas --- Reporte de Desviaciones:}
{\tiny
\renewcommand{\arraystretch}{1.15}
\begin{longtable}{|L{1.1cm}|L{2.3cm}|L{3.2cm}|L{3.6cm}|L{1.2cm}|L{3.6cm}|}
\hline
\textbf{ID} & \textbf{Caso} & \textbf{Descripción} & \textbf{Resultado esperado} & \textbf{Aprobado} & \textbf{Observaciones} \\ \hline
\endhead
DVS-01 & Apertura & Seleccionar «Desviaciones». & Se abre el modal clasificado. & Sí & \mth{test_dvs01_apertura}: se abrió el Toplevel «Reporte de desviaciones» de 780x560 con el subtítulo «Desviaciones de tiempo por diapositiva» y las dos columnas. \\ \hline
DVS-02 & Duración insuficiente & Usar tiempos reales menores. & Se muestran los tres déficits principales. & Sí & \mth{test_dvs02_duracion_insuficiente}: la diapositiva 1 (real 20s, IA 60s) apareció con «-67\,\%». \\ \hline
DVS-03 & Duración excesiva & Usar tiempos reales mayores. & Se muestran los tres excesos principales. & Sí & \mth{test_dvs03_duracion_excesiva}: la diapositiva 2 (real 200s, IA 120s) apareció con «+67\,\%». \\ \hline
DVS-04 & Porcentaje & Comparar con un cálculo manual. & Los porcentajes coinciden. & Sí & \mth{test_dvs04_porcentaje_correcto}: la diapositiva 4 (real 40s, IA 150s) mostró «-73\,\%», igual al cálculo manual. \\ \hline
DVS-05 & Filtro & Incluir diapositivas no presentadas. & No se clasifican. & Sí & \mth{test_dvs05_diapositivas_no_presentadas_excluidas}: la diapositiva con real 0 no apareció en ninguna columna. \\ \hline
DVS-06 & Categoría vacía & Usar un grupo sin déficit. & Se informa la ausencia de casos. & Sí & \mth{test_dvs06_categoria_vacia}: la columna sin datos mostró «Sin diapositivas en esta categoría.». \\ \hline
DVS-07 & Cierre con Entendido & Seleccionar el botón inferior. & Se cierra solo el modal. & Sí & \mth{test_dvs07_cierre_con_entendido}: al invocar «Entendido» el modal dejó de existir. \\ \hline
DVS-08 & Cierre con X & Seleccionar la X. & Se cierra solo el modal. & Sí & \mth{test_dvs08_cierre_con_la_x}: la rutina de \texttt{WM\_DELETE\_WINDOW} cerró el modal. \\ \hline
DVS-09 & Reapertura & Cerrar y abrir nuevamente. & Se reconstruye sin residuos. & Sí & \mth{test_dvs09_reapertura}: la segunda apertura creó un modal nuevo y válido. \\ \hline
DVS-10 & Cambio de grupo con el modal abierto & Con el modal abierto, cambiar de grupo. & Se muestran los datos del nuevo grupo. & No & \mth{test_usr01_clic_repetido_en_desviaciones} y \mth{test_usr09_salir_con_modal_abierto} confirmaron que el modal retiene el \texttt{grab}, por lo que las pestañas quedan bloqueadas hasta cerrarlo. \\ \hline
DVS-11 & Cambio de grupo previo & Cambiar de grupo y abrir «Desviaciones». & Se muestran los datos del grupo activo. & Sí & \mth{test_dvs11_cambio_de_grupo_previo}: el encabezado cambió de «REPORTE · 3CM2» a «REPORTE · 3CM4». \\ \hline
DVS-12 & Límite de tres por categoría & Usar más de tres desviaciones por categoría. & Solo se listan las tres más extremas. & Sí & \mth{test_dvs12_limite_de_tres_por_categoria}: con cinco déficits solo se listaron tres diapositivas. \\ \hline
DVS-13 & Referencia IA igual a cero & Incluir una diapositiva sin tiempo recomendado. & La diapositiva no se clasifica. & Sí & \mth{test_dvs13_referencia_ia_cero}: la diapositiva con IA = 0 no generó desviación porcentual ni apareció en las columnas. \\ \hline
DVS-14 & Encabezado y detalle & Revisar el encabezado y las filas. & Indica el grupo y el detalle. & Sí & \mth{test_dvs14_encabezado_y_detalle}: se verificó «REPORTE · 3CM2», el detalle «Real … · IA …» y la zona del análisis («Introduccion»). \\ \hline
DVS-15 & Sin grupo activo & Abrir el modal sin grupo. & Se informa la ausencia de datos. & Sí & \mth{test_dvs15_sin_grupo}: con grupo nulo el encabezado mostró «REPORTE · —». \\ \hline
\caption{Cobertura de pruebas del módulo de reporte de desviaciones.}
\label{tab:cobertura-desviaciones}
\end{longtable}
}

\clearpage
\textbf{Módulo: Pruebas de Integridad y Usuario}

\textbf{Objetivo:} Validar la consistencia entre ambos JSON y comprobar que la interacción repetida no produzca estados inconsistentes.

\textbf{Descripción del flujo:}
\begin{enumerate}
\item Verificar los índices en base uno y la correspondencia entre JSON.
\item Comparar los tiempos individuales con el total.
\item Ejecutar clics, desplazamientos y cambios de grupo repetidos.
\end{enumerate}

\textbf{Componentes que probar:}
\begin{table}[H]
\centering
\begin{tabular}{|L{0.34\textwidth}|L{0.18\textwidth}|L{0.38\textwidth}|}
\hline
\textbf{Componente} & \textbf{Tipo} & \textbf{Descripción} \\ \hline
Índices de diapositiva & Datos & Verifica el uso de base uno. \\ \hline
Correspondencia entre JSON & Datos & Compara el análisis con los tiempos registrados. \\ \hline
Interfaz de usuario & UI/Frontend & Evalúa eventos rápidos y repetidos. \\ \hline
\end{tabular}
\caption{Componentes involucrados en las pruebas de integridad y usuario del prototipo 3.}
\label{tab:componentes-integridad-tt3}
\end{table}

\textbf{Cobertura de pruebas --- Integridad y Usuario:}
{\tiny
\renewcommand{\arraystretch}{1.15}
\begin{longtable}{|L{1.1cm}|L{2.3cm}|L{3.2cm}|L{3.6cm}|L{1.2cm}|L{3.6cm}|}
\hline
\textbf{ID} & \textbf{Caso} & \textbf{Descripción} & \textbf{Resultado esperado} & \textbf{Aprobado} & \textbf{Observaciones} \\ \hline
\endhead
INT-01 & Base uno & Revisar las claves de tiempos. & La clave 1 representa la primera diapositiva. & Sí & \mth{test_int01_base_uno}: las diapositivas indexan 1..5 y las claves del reporte son «1»..«5». \\ \hline
INT-02 & Coincidencia & Comparar ambos JSON. & La gráfica no presenta huecos inesperados. & Sí & \mth{test_int02_correspondencia_entre_json}: el punto 2 tomó el real del grupo (200 s) y la IA del análisis (120 s). \\ \hline
INT-03 & Suma total & Sumar los tiempos individuales. & La suma coincide con la duración real. & Sí & \mth{test_int03_suma_total}: al faltar el total, el dashboard sumó 30+40 = 70 s y mostró «1m 10s». \\ \hline
INT-04 & Diapositiva omitida & Consultar una diapositiva omitida. & Se identifica como no presentada. & Sí & \mth{test_int04_diapositiva_omitida}: el detalle de la diapositiva 5 incluyó «(no presentada)». \\ \hline
INT-05 & Pestañas & Alternar diez veces entre grupos. & No se acumulan componentes. & Sí & \mth{test_int05_pestanas_estables}: tras alternar 10 veces el número de pestañas permaneció en 2. \\ \hline
INT-06 & Zona por diapositiva & Revisar la zona en las tarjetas de desviación. & Proviene del análisis. & Sí & \mth{test_int06_zona_por_diapositiva}: la fila mostró la zona «Desarrollo» tomada de \texttt{zona\_slide}. \\ \hline
INT-07 & Total derivado & Abrir un análisis sin \texttt{total\_diapositivas}. & El dashboard calcula el total. & Sí & \mth{test_int07_total_derivado}: sin el campo total, el dashboard derivó 3 del máximo \texttt{slide\_number}. \\ \hline
INT-08 & Tiempo IA derivado & Abrir un análisis sin el tiempo total. & Se suma por diapositiva. & Sí & \mth{test_int08_tiempo_ia_derivado}: sin \texttt{tiempo\_total\_exposicion\_segundos} el tiempo IA se calculó como 10+20+30 = 60 s («1m 00s»). \\ \hline
USR-01 & Clic repetido en Desviaciones & Realizar varios clics rápidos. & Se controla la apertura del modal. & Sí & \mth{test_usr01_clic_repetido_en_desviaciones}: mientras el modal existe, \texttt{grab\_current()} apunta a él, de modo que la gráfica no recibe más clics. \\ \hline
USR-02 & Clic repetido en Volver & Pulsar varias veces «Volver a inicio». & La transición se ejecuta una vez. & Sí & \mth{test_usr02_clic_repetido_en_volver}: la segunda llamada no produjo error y la vista quedó en «main». \\ \hline
USR-03 & Desplazamiento veloz & Mover rápidamente la barra. & La gráfica responde sin errores. & Sí & \mth{test_usr03_desplazamiento_veloz}: seis saltos de ventana consecutivos no generaron errores. \\ \hline
USR-04 & Hover prolongado & Mantener el cursor sobre un punto. & El detalle permanece estable. & Sí & \mth{test_usr04_hover_repetido}: seis eventos de movimiento consecutivos no acumularon elementos «hover». \\ \hline
USR-05 & Grupo y hover & Cambiar de grupo y mover el cursor. & El detalle usa el grupo actual. & Sí & \mth{test_usr05_grupo_y_hover}: tras cambiar a \texttt{3CM4} el detalle y las etiquetas mostraron el valor del nuevo grupo («1m39s» con 99 s). \\ \hline
USR-06 & Redimensionado y hover & Redimensionar con el detalle visible. & El detalle se limpia o reposiciona. & Sí & \mth{test_usr06_redimensionado_con_hover}: el redimensionado con el detalle visible reprogramó el dibujado sin errores. \\ \hline
USR-07 & Grupo sin excesos & Abrir «Desviaciones» con un grupo sin excesos. & Una columna queda con mensaje vacío. & Sí & \mth{test_usr07_grupo_sin_excesos}: la columna de excesos mostró «Sin diapositivas en esta categoría.». \\ \hline
USR-08 & Redimensionado repetido & Redimensionar varias veces seguidas. & No se acumulan redibujados. & Sí & \mth{test_usr08_redimensionado_repetido}: varios \texttt{<Configure>} seguidos dejaron un único redibujado pendiente. \\ \hline
USR-09 & Salir con el modal abierto & Regresar al inicio con el modal abierto. & La vista se restaura sin residuos. & Sí & \mth{test_usr09_salir_con_modal_abierto}: el modal retiene el \texttt{grab} (la topbar no responde hasta cerrarlo); tras cerrarlo, el regreso a «main» se completó sin residuos. \\ \hline
\caption{Cobertura de pruebas de integridad y usuario del prototipo 3.}
\label{tab:cobertura-integridad-tt3}
\end{longtable}
}

% ===========================================================================
\clearpage
\section*{Trazabilidad de las pruebas}
\addcontentsline{toc}{section}{Trazabilidad de las pruebas}

La Tabla \ref{tab:trazabilidad-tt3} vincula cada identificador del protocolo con su prueba automatizada. Todas se ejecutan con \texttt{python -m tests.e2e.run\_e2e} y su salida queda en \texttt{docs/prototipo\_final/evidencias/}.

{\tiny
\renewcommand{\arraystretch}{1.1}
\begin{longtable}{|L{1.4cm}|L{9.5cm}|L{2.2cm}|}
\hline
\textbf{ID} & \textbf{Prueba automatizada} & \textbf{Resultado} \\ \hline
\endhead
PRF-01 & \mth{test_prf01_entrada_a_la_vista} & Pasa \\ \hline
PRF-02 & \mth{test_prf02_carga_de_analisis} & Pasa \\ \hline
PRF-03 & \mth{test_prf03_carga_de_tiempos} & Pasa \\ \hline
PRF-04 & \mth{test_prf04_resumen_con_datos} & Pasa \\ \hline
PRF-05 & \mth{test_prf05_resumen_sin_grupos} & Pasa (panel) \\ \hline
PRF-06 & \mth{test_prf06_sin_analisis} & Pasa (panel) \\ \hline
PRF-07 & \mth{test_prf07_tiempo_ia_es_la_suma} & Pasa \\ \hline
PRF-08 & \mth{test_prf08_grupo_inicial} & Pasa \\ \hline
PRF-09 & \mth{test_prf09_analisis_corrupto} & Pasa \\ \hline
PRF-10 & \mth{test_prf10_reporte_corrupto} & Pasa \\ \hline
PRF-11 & \mth{test_prf11_regreso_al_inicio} & Pasa \\ \hline
PRF-12 & \mth{test_prf12_eliminacion_no_aplica} & Pasa (N/A) \\ \hline
PRF-13 & \mth{test_prf13_calificacion_por_etiqueta} & Pasa \\ \hline
PRF-14 & \mth{test_prf14_desviacion_global_por_umbral} & Pasa \\ \hline
PRF-15 & \mth{test_prf15_diapositivas_presentadas} & Pasa \\ \hline
PRF-16 & \mth{test_prf16_planeado_usa_la_ultima_sesion} & Pasa \\ \hline
PRF-17 & \mth{test_prf17_reapertura_sin_residuos} & Pasa \\ \hline
GRF-01 & \mth{test_grf01_dibujado_inicial} & Pasa \\ \hline
GRF-02 & \mth{test_grf02_area_real} & Pasa \\ \hline
GRF-03 & \mth{test_grf03_serie_ia_discontinua} & Pasa \\ \hline
GRF-04 & \mth{test_grf04_tooltip_muestra_detalle} & Pasa \\ \hline
GRF-05 & \mth{test_grf05_cursor_lejos_del_punto} & Pasa \\ \hline
GRF-06 & \mth{test_grf06_desplazamiento_entre_ventanas} & Pasa \\ \hline
GRF-07 & \mth{test_grf07_cambio_de_grupo_redibuja} & Pasa \\ \hline
GRF-08 & \mth{test_grf08_sincronizacion_con_el_resumen} & Pasa \\ \hline
GRF-09 & \mth{test_grf09_escala_vertical} & Pasa \\ \hline
GRF-10 & \mth{test_grf10_sin_grupos} & Pasa (N/A) \\ \hline
GRF-11 & \mth{test_grf11_redimensionado} & Pasa \\ \hline
GRF-12 & \mth{test_grf12_muchas_diapositivas} & Pasa \\ \hline
GRF-13 & \mth{test_grf13_diapositiva_no_presentada} & Pasa \\ \hline
GRF-14 & \mth{test_grf14_tiempo_ia_cero_no_rompe} & Pasa \\ \hline
GRF-15 & \mth{test_grf15_etiquetas_de_valor} & Pasa \\ \hline
GRF-16 & \mth{test_grf16_leyenda} & Pasa \\ \hline
GRF-17 & \mth{test_grf17_limites_del_desplazamiento} & Pasa \\ \hline
GRF-18 & \mth{test_grf18_presentacion_corta} & Pasa \\ \hline
DVS-01 & \mth{test_dvs01_apertura} & Pasa \\ \hline
DVS-02 & \mth{test_dvs02_duracion_insuficiente} & Pasa \\ \hline
DVS-03 & \mth{test_dvs03_duracion_excesiva} & Pasa \\ \hline
DVS-04 & \mth{test_dvs04_porcentaje_correcto} & Pasa \\ \hline
DVS-05 & \mth{test_dvs05_diapositivas_no_presentadas_excluidas} & Pasa \\ \hline
DVS-06 & \mth{test_dvs06_categoria_vacia} & Pasa \\ \hline
DVS-07 & \mth{test_dvs07_cierre_con_entendido} & Pasa \\ \hline
DVS-08 & \mth{test_dvs08_cierre_con_la_x} & Pasa \\ \hline
DVS-09 & \mth{test_dvs09_reapertura} & Pasa \\ \hline
DVS-10 & \mth{test_usr01_clic_repetido_en_desviaciones} & Pasa (documenta el bloqueo) \\ \hline
DVS-11 & \mth{test_dvs11_cambio_de_grupo_previo} & Pasa \\ \hline
DVS-12 & \mth{test_dvs12_limite_de_tres_por_categoria} & Pasa \\ \hline
DVS-13 & \mth{test_dvs13_referencia_ia_cero} & Pasa \\ \hline
DVS-14 & \mth{test_dvs14_encabezado_y_detalle} & Pasa \\ \hline
DVS-15 & \mth{test_dvs15_sin_grupo} & Pasa \\ \hline
INT-01 & \mth{test_int01_base_uno} & Pasa \\ \hline
INT-02 & \mth{test_int02_correspondencia_entre_json} & Pasa \\ \hline
INT-03 & \mth{test_int03_suma_total} & Pasa \\ \hline
INT-04 & \mth{test_int04_diapositiva_omitida} & Pasa \\ \hline
INT-05 & \mth{test_int05_pestanas_estables} & Pasa \\ \hline
INT-06 & \mth{test_int06_zona_por_diapositiva} & Pasa \\ \hline
INT-07 & \mth{test_int07_total_derivado} & Pasa \\ \hline
INT-08 & \mth{test_int08_tiempo_ia_derivado} & Pasa \\ \hline
USR-01 & \mth{test_usr01_clic_repetido_en_desviaciones} & Pasa \\ \hline
USR-02 & \mth{test_usr02_clic_repetido_en_volver} & Pasa \\ \hline
USR-03 & \mth{test_usr03_desplazamiento_veloz} & Pasa \\ \hline
USR-04 & \mth{test_usr04_hover_repetido} & Pasa \\ \hline
USR-05 & \mth{test_usr05_grupo_y_hover} & Pasa \\ \hline
USR-06 & \mth{test_usr06_redimensionado_con_hover} & Pasa \\ \hline
USR-07 & \mth{test_usr07_grupo_sin_excesos} & Pasa \\ \hline
USR-08 & \mth{test_usr08_redimensionado_repetido} & Pasa \\ \hline
USR-09 & \mth{test_usr09_salir_con_modal_abierto} & Pasa \\ \hline
\caption{Trazabilidad entre los identificadores del protocolo y las pruebas automatizadas.}
\label{tab:trazabilidad-tt3}
\end{longtable}
}

% ===========================================================================
\clearpage
\section*{Notas de la ejecución}
\addcontentsline{toc}{section}{Notas de la ejecución}

\begin{itemize}
    \item La corrida documentada aquí es la de \texttt{docs/prototipo\_final/evidencias/20261018\_233448/}: 97 pruebas, 0 fallos, 0 errores, resultado OK (15 de negocio + 15 de interfaz + 67 de rendimiento).
    \item Las pruebas PRF-05, PRF-06 y GRF-10 se marcan como \emph{No} o \emph{N/A} en su resultado funcional porque el control de acceso impide abrir el dashboard sin sesiones (por diseño). Su comportamiento degradado sí quedó validado a nivel de panel por \mth{test_prf05_resumen_sin_grupos}, \mth{test_prf06_sin_analisis} y \mth{test_grf10_sin_grupos}.
    \item DVS-10 se marca \emph{No} porque el modal usa \texttt{grab\_set}: para cambiar de grupo es necesario cerrar antes el reporte de desviaciones. Es un comportamiento deliberado de la ventana modal y quedó evidenciado por \mth{test_usr01_clic_repetido_en_desviaciones} y \mth{test_usr09_salir_con_modal_abierto}.
    \item Los tiempos IA provienen del JSON de análisis y los reales del reporte de tiempos; cuando un campo total falta, el dashboard los deriva por suma (INT-03, INT-07, INT-08).
    \item La vista de rendimiento activa es \texttt{performance\_gui.py}; el archivo \texttt{rendimiento\_gui.py} del repositorio corresponde a una variante anterior que no se utiliza.
\end{itemize}

\end{document}
